\documentclass[11pt,a4paper]{article}
\usepackage[T1]{fontenc}
\usepackage{lmodern,amsmath,amssymb,amsthm,mathtools,booktabs,microtype}
\usepackage[margin=25mm]{geometry}
\usepackage[colorlinks=true,linkcolor=blue,urlcolor=blue,citecolor=blue]{hyperref}
\newtheorem{theorem}{Theorem}[section]
\newtheorem{lemma}[theorem]{Lemma}
\newtheorem{proposition}[theorem]{Proposition}
\newcommand{\Q}{\mathbb Q}\newcommand{\Z}{\mathbb Z}
\setlength{\emergencystretch}{2em}
\hypersetup{pdftitle={Smaller Coefficients for a Quartic in 214 Variables}}
\title{Smaller Coefficients for a Quartic in 214 Variables}
\author{}\date{13 November 2026}
\begin{document}\maketitle
\begin{abstract}
An invertible rational rescaling of six witness coordinates reduces the
coefficient growth of the specified 214-variable quartic family. The
original coefficient height is $149^{4N+2}+1$ for $N\ge53$. The new height
is exactly $31581660^{28}149^{2N}$ for $N\ge57$ and is strictly smaller
than the original height for every $N\ge53$. The rational input projection
is unchanged for each parameter. All 188 quadratic equations and all
2919 quartic coefficient profiles are checked exactly. The argument is
algebraic: it does not verify the original arithmetic foundations or
make the sufficient defining threshold effective.
\end{abstract}

\section{The source and the statement}
The source polynomial is the sum of squares of 188 primitive integer
quadratic equations in 214 rational coordinates. Its first coordinate
is the distinguished input. Write the fixed coefficient as
\[
 A_N=149^{2N+1},\qquad
 P_N(X_1,\ldots,X_{214})=\sum_{i=1}^{188}r_i(A_N,\boldsymbol X)^2.
\]
The exact source version is the family retained in the complete
Version 17 manuscript at commit
\begin{center}\small
\nolinkurl{1f7e1f1fb0438feda3c9f3da9cc16982126aeeed}.
\end{center}
Its published coefficient file has Git blob identifier
\begin{center}\small
\nolinkurl{f6a32cc2f7ebcc37a952c5ce737169608cfe27c2}.
\end{center}
The accompanying source table is reconstructed from the printed equations.
Its sum of squares reproduces that exact published coefficient file;
no assertion of byte identity with the source's separate quadratic table
is needed. The source and all new tables are retained in
\begin{center}\small
\nolinkurl{computation/quartic/balanced214_v24/}.
\end{center}

Set
\[
 K=31581660^{28},\qquad J=5898050446009651968,\qquad H_0=(JK)^2.
\]
Coefficient height means the largest absolute integer coefficient after
substitution of the fixed parameter. All coordinate counts include the input.

\begin{theorem}\label{thm:main}
There is a completely specified integer quartic family
$\widetilde P_N\in\Z[Y_1,\ldots,Y_{214}]$, for integers $N\ge1$, with
these properties.
\begin{enumerate}
\item For every $x\in\Q$ and every $N\ge1$,
\[
 \exists\boldsymbol y\in\Q^{213}\ P_N(x,\boldsymbol y)=0
 \quad\Longleftrightarrow\quad
 \exists\boldsymbol z\in\Q^{213}\ \widetilde P_N(x,\boldsymbol z)=0.
\]
\item Each member has degree exactly four, integer content one, and exactly
2919 nonzero coordinate monomials. It is a sum of 188 primitive quadratic
squares.
\item For every $N\ge1$,
\[
 H(\widetilde P_N)=\max\{H_0,K149^{2N}\}.
\]
For $N\ge57$ this is $K149^{2N}$. For every $N\ge53$,
\[
 H(\widetilde P_N)<149^{4N+2}+1=H(P_N).
\]
\end{enumerate}
\end{theorem}
The sufficient defining threshold of the source is not evaluated here.
Consequently, if its integer-definability theorem is used with its stated
premises and a fixed ineffective threshold $N_0$, exactly the same
$N_0$ works for the new family. Neither $53$ nor $57$ is asserted to be
a sufficient defining threshold. This is a coefficient improvement of
that source, not a new variable-count result or an independent proof of
its global arithmetic theorem.

\section{Rational rescaling and row normalization}
Put $C=149^N$, so $A_N=149C^2$, and let
\[
 I=\{90,91,93,117,118,120\}.
\]
Make the coordinate substitution
\begin{equation}\label{eq:substitution}
 X_j=\begin{cases}CY_j,&j\in I,\\Y_j,&j\notin I.\end{cases}
\end{equation}
This is an invertible rational linear map for every nonzero rational $C$;
it fixes $X_1=Y_1$.

For the source row numbers, define
\[
 m_{35}=2,\qquad m_{36}=m_{37}=m_{38}=m_{60}=m_{63}=m_{65}=1,
\]
and $m_i=0$ for all other rows. The new primitive rows are
\begin{equation}\label{eq:rows}
 \widetilde r_i(C,\boldsymbol Y)
 =\varepsilon_i C^{-m_i}
 r_i(149C^2,\boldsymbol X(C,\boldsymbol Y)),
 \qquad \varepsilon_i\in\{1,-1\}.
\end{equation}
The sign is chosen so the first coefficient in the prescribed monomial
order is positive; it has no effect on the sum of squares. Each numerator
in~\eqref{eq:rows} is divisible by the stated power of $C$ in
$\Z[C,\boldsymbol Y]$. The quotient has integer content one, degree at
most two in $\boldsymbol Y$, and degree at most one in $C$.

These statements can be checked term by term without evaluating any large
parameter. A source term $aA^eX_{j_1}\cdots X_{j_k}$ becomes
\[
 a149^e C^{2e+\#\{\ell:j_\ell\in I\}}
 Y_{j_1}\cdots Y_{j_k}.
\]
For each of the 188 rows the minimum exponent is $m_i$, and the difference
between maximum and minimum exponents is at most one. The integer
coefficients after division have greatest common divisor one. The supplied
checker verifies the displayed identity for every term and every row.
No coordinate monomial in an individual row has both a constant and a
linear coefficient in $C$. If $g_0,g_1$ are the greatest common divisors
of the constant and linear coefficient groups of a row, the checker gives
$\gcd(g_0,g_1)=\gcd(g_0,149)=1$. Thus every specialization $C=149^N$
remains primitive over $\Z$. Fifty-two rows change; the remaining 136
retain their coefficient vectors up to the harmless chosen signs.

Define
\[
 \widetilde P_N(\boldsymbol Y)=
 \sum_{i=1}^{188}\widetilde r_i(149^N,\boldsymbol Y)^2.
\]
Over $\Q$, a sum of squares vanishes exactly when each summand vanishes.
All the row multipliers in~\eqref{eq:rows} are nonzero, and the coordinate
map~\eqref{eq:substitution} is bijective. This proves both directions of
the projection assertion, on the entire rational solution set. No witness
coordinate is divided out; zero witness values remain allowed. No local
or global arithmetic existence assertion is added.

\section{Exact coefficient height}
Collect the expansion by its coordinate monomials $M$. Write
\[
 \widetilde P(C,\boldsymbol Y)=\sum_M p_M(C)M(\boldsymbol Y).
\]
The exact expansion has 2924 terms before equal coordinate monomials are
collected, and 2919 after collection. The coefficient types are
\begin{center}
\begin{tabular}{cr}
\toprule
Nonzero powers of $C$ in $p_M$ & Number of coordinate monomials\\\midrule
$\{0\}$&2508\\
$\{1\}$&331\\
$\{2\}$&75\\
$\{0,2\}$&5\\\bottomrule
\end{tabular}
\end{center}
In each of the five mixed profiles the two coefficients have the same
sign. Therefore no coordinate monomial disappears when $C>0$.
Every one of the 214 coordinates occurs. The coefficient of $Y_1^4$
is exactly one for every $C$, proving degree four and content one.

The two coefficients
\begin{equation}\label{eq:attaining}
 p_{Y_{112}^4}(C)=H_0,\qquad
 p_{Y_{109}^2Y_{117}^2}(C)=KC^2
\end{equation}
will attain the height. The following elementary observation makes the
check for all parameters finite.

\begin{lemma}\label{lem:envelope}
Let $H_0,K>0$ and $C>0$.
\begin{enumerate}
\item If $a^2\le H_0K$, then $|a|C\le\max(H_0,KC^2)$.
\item If $|a_0|K+|a_2|H_0\le H_0K$, then
$|a_0+a_2C^2|\le\max(H_0,KC^2)$.
\end{enumerate}
\end{lemma}
\begin{proof}
The geometric mean of $H_0$ and $KC^2$ is no larger than their maximum,
which proves the first assertion. For the second, put $u=|a_0|/H_0$
and $v=|a_2|/K$. Then $u,v\ge0$ and $u+v\le1$. The triangle
inequality gives
\[
 |a_0+a_2C^2|\le uH_0+vKC^2\le\max(H_0,KC^2).
\]
\end{proof}

For each of the 331 linear profiles the exact integer calculation gives
$a^2\le H_0K$. For each of the remaining 2588 profiles it gives
$|a_0|K+|a_2|H_0\le H_0K$, with a missing coefficient taken to be zero.
Every coefficient is therefore bounded by the claimed maximum, and
\eqref{eq:attaining} proves equality. The verification concerns every
coefficient, not just an upper bound from a circuit norm.

To locate the integer threshold, the checker verifies
\[
 149^{56}<J31581660^{14}<149^{57}.
\]
Since $H_0/K=J^2K$, it follows that $KC^2>H_0$ for every
$C=149^N$ with $N\ge57$, and that $H_0>KC^2$ for $N\le56$.
The strict improvement over the source follows from the two exact
inequalities
\[
 H_0<149^{214}+1,\qquad K<149^{108}.
\]
For $N\ge53$, the first is bounded by $149^{4N+2}+1$, and the second
gives $K149^{2N}<149^{4N+2}$. These prove the height assertions of
Theorem~\ref{thm:main}.

For completeness the source height is also rechecked: its unique
quadratic profile in $A$ is $A^2+1$. Every other profile has degree at
most one and absolute value strictly smaller at $A=149^{107}$.
After division by $A^2$, the corresponding absolute-value bounds decrease
as $A$ increases. This establishes $H(P_N)=149^{4N+2}+1$ for all
$N\ge53$, independently of the new coefficient calculation.

\section{Reproducibility and audit scope}
The accompanying standard-library producer writes the full new quadratic
and quartic tables. A separate checker imports no producer and expands
squares using unordered term pairs rather than the producer's ordered
double loop. It performs the following finite checks: source polynomial
identity, expansion of all reconstructed source rows, every row transformation,
all new quartic coefficients, the exact support and content, all height
envelope inequalities, and both parameter thresholds. A deliberately
corrupted transformed coefficient is rejected.

The reconstructed source polynomial has 246349 bytes and SHA-256
\begin{center}\small
\nolinkurl{9f461531df52ac73e02cf7a4ce270422715642c0d3112952cad4968d9a7e24ea}.
\end{center}
Its Git blob identifier matches the live repository metadata at the pinned
commit. This binds the numerical source used here to the actual published
polynomial, not merely to a count or an old review verdict. Coordinate
names and the reconstructed equations are distributed beside the checker.
The named indices in the proof are the original source indices, so no
coordinate renumbering is implicit.

These algebraic checks do not establish the source's elliptic rank,
prime-pattern, pole-realization or contact-height theorems. In particular
they do not supply complete independent acceptance of the assembled
arithmetic proof. The producer, checker and written argument were developed
in one research session; distinct implementations are not presented as
non-author review. The statement proved here is exact algebraic
projection equivalence and a coefficient improvement, with the original
arithmetic status and quantifier order left unchanged.
\end{document}
